\chapter{Grundlagen}\label{ch:grundlagen}

%TODO text einführung

\section{Statische und interaktive 3D-Objekte}\label{sec:statische_interaktive_3d_objekte}

Im Folgenden werden statische sowie interaktive 3D-Objekte definiert und erklärt.
Die Begriffe 3D-Objekt und 3D-Asset werden in dieser Arbeit simultan verwendet. Beide beschreiben ein Element einer virtuellen Szene und können als alleinige Geometrie existieren oder Unterkomponenten bzw. Teilobjekte besitzen.

\subsection*{Statische 3D-Objekte}
In dieser Arbeit werden statische 3D-Objekte als Basisform digitaler Modelle betrachtet.
Sie sind die fundamentale Basis, auf denen komplexere Verhaltensweisen aufbauen können, jedoch wohnt ihnen mit diesen Informationen noch kein Eigenleben inne.
Statische 3D-Objekte lassen sich meist rein geometrisch beschreiben.
Die Mesh-Geometrie besteht aus Dreiecksnetzen oder Polyeder-Darstellungen, welche die Oberfläche durch polygonale Flächen, in der Regel Dreiecksnetze, deren Eckpunkte (\textit{Vertices}) und Verbindungen (\textit{Edges/Faces}) definieren~\cite{botschPolygonMeshProcessing2010}.
Weiterhin besitzen sie grafische Attribute wie Materialeigenschaften, Texturen, Transparenz und Sichtbarkeit sowie eine hierarchische Anordnung im Szenengraph, welcher ihre Position, Rotation und Skalierung organisiert~\cite{GabrielZachmannsDissertation}.

\subsection*{Interaktive 3D-Objekte}

Während statische 3D-Objekte als unbelebte geometrische Körper betrachtet werden können, zeichnen sich interaktive 3D-Objekte durch Wissen über ihre Funktionalität und die Möglichkeit zur Interaktion mit ihnen aus~\cite{kallmannModelingBehaviorsInteractive2002}.
Auf diese 3D-Objekte kann durch Eingaben (\textit{Input}) eingewirkt werden und das Objekt reagiert zustands- oder verhaltensbezogen~\cite{steuerDefiningVirtualReality1992}.
Durch diese Reaktion wird aus dem statischen 3D-Objekt ein interaktionsfähiges Element.
Es ist somit in der Lage, auf durch Benutzereingaben ausgelöste Ereignisse (\textit{Events}) zu reagieren und wird zu einem handlungsoffenen Teil der virtuellen Welt. Interaktive 3D-Objekte werden in dieser Arbeit auch als interaktive Prototypen bezeichnet. %TODO auf duch Benutzer*innen

\section{Transformation statischer in interaktive 3D-Objekte}

Die in der vorliegenden Arbeit untersuchte Transformation von statisch zu interaktiv ist keine rein geometrische oder visuelle Modifikation, sondern eine semantisch-verhaltensbezogene.
Sie macht aus einem visuellen und geometrischen Asset ein interaktives Modell, welches Informationen über sein eigenes Verhalten und die möglichen Interaktionen mit diesem beinhaltet.
Mit dem Konzept dieser \textit{Smart Objects} wird Logik direkt im Objekt gespeichert, was wiederum Dezentralisierung und Wiederverwendbarkeit ermöglicht~\cite{kallmannModelingBehaviorsInteractive2002}. Dieser Prozess umfasst verschiedene Ebenen.

Zunächst wird das statische 3D-Modell segmentiert sowie semantisch etikettiert, was sich mit dem Begriff \textit{Labeling} zusammenfassen lässt.
Dabei werden rohe Mesh-Geometrien in bedeutungsvolle Teile mit bestimmten Rollen zerlegt (\zB~der Griff einer Tasse oder die Sitzfläche eines Stuhls), sodass ein System später \enquote{versteht}, welches Teil eines 3D-Objekts für welche Interaktion relevant ist und eine Zuordnung stattfinden kann~\cite{moPartNetLargescaleBenchmark2018}\cite{hassaninVisualAffordanceFunction2021}.

Im nächsten Schritt erfolgt der Übergang von der Semantik zum Verhalten mittels Zuweisung spezifischer Interaktionsmerkmale (\textit{Interaction Features}) zu den im ersten Schritt analysierten semantischen Etiketten. Diese Merkmale beinhalten eine Beschreibung der Funktionalität sowie eine Verknüpfung mit Verhaltensplänen oder Zustandsautomaten. Diese Pläne können Zustände überprüfen oder ändern und legen das erwartete Verhalten von beteiligten Akteuren fest~\cite{kallmannModelingBehaviorsInteractive2002}.

\section{Das Vivian Framework}\label{sec:vivian_framework}

Das Vivian Framework ist ein \textit{Virtual-Prototyping-Framework}.
Es erstellt virtuelle Prototypen, also digitale, interaktive Repräsentationen physischer Objekte aus der realen Welt.
Diese Repräsentation kann \zB~in der Produktentwicklung verwendet werden, um virtuelle Prototypen von echten Produkten zu erzeugen und deren Verhalten frühzeitig zu testen, ohne dass das physische Produkt tatsächlich vorhanden sein muss.
Diese virtuellen Prototypen funktionieren spezifikationsgetrieben, sodass Verhalten nicht programmatisch, sondern deklarativ als Spezifikation geschrieben wird.
Zur Laufzeit wird diese Spezifikation dann mittels einer \textit{State-Machine}-basierten Umgebung interpretiert, um ein Interaktionsverhalten abzubilden~\cite{VivianFrameworkGitLab2021}.
Ein Vorteil einer solchen spezifikationsgetriebenen Vorgehensweise ist, dass beispielsweise Produktentwickler*innen keine Programmierkenntnisse benötigen, um virtuelle Prototypen von Produkten (\textit{Digital Twins}) zu erstellen.
So kann das Verhalten eines Produktes bereits frühzeitig und iterativ getestet werden.
Ein weiteres mögliches Einsatzgebiet ist das Usability Engineering, mit dem interaktive Systeme so gestaltet und geprüft werden, dass sie für die jeweilige Zielgruppe möglichst effektiv, effizient sowie zufriedenstellend nutzbar sind.
Dafür werden Systeme nicht erst am Ende, sondern bereits in der Entwicklung, beispielsweise mittels Prototypen, hinsichtlich Gebrauchstauglichkeit analysiert und getestet~\cite{gouldDesigningUsabilityKey1985}.

\subsection*{Vivian-Spezifikationen}

Wie bereits gezeigt, besitzen statische 3D-Objekte von sich aus keine Interaktivität (siehe~\ref{sec:statische_interaktive_3d_objekte}).
Das Vivian-Framework bildet die Brücke zwischen statischem und interaktivem 3D-Objekt über eine maschinenlesbare Beschreibung des Verhaltens bei Interaktion mit dem virtuellen Prototyp.
Die Verbindung dieser Spezifikation zu einem 3D-Objekt in der Szene wird über Namensreferenzen hergestellt. %TODO vivian framework gleiche schreibweise überall?
Ein Element in der Spezifikation referenziert demnach ein gleichnamiges 3D-Objekt oder Teilobjekt in der Szene.
Weiterhin ermöglichen Spezifikationen eine Entkopplung des Verhaltensmodells vom 3D-Objekt, sodass unterschiedliche Spezifikationen verschiedene Verhalten bei demselben 3D-Objekt ermöglichen.
Sie besteht im Wesentlichen aus vier Teilen: Interaktionselemente (\textit{Interaction Elements}), Visualisierungselemente (\textit{Visualization Elements}), Zustände (\textit{States}) und Übergänge (\textit{Transitions})~\cite{VivianFrameworkGitLab2021}.

%TODO texttt weg für elemente
\textbf{Interaktionselemente.} Hier befinden sich Elemente, welche Eingaben akzeptieren und somit Interaktion ermöglichen. Das können \ua \texttt{Buttons}, \texttt{Slider} oder \texttt{Movables} sein.

\textbf{Visualisierungselemente.} Dabei handelt es sich um Elemente, welche eine visuelle Rückmeldung geben. Typisch sind hier \texttt{Light}, \texttt{Screen} oder \texttt{Sound Source}.

\textbf{Zustände.} In diesem Teil befinden sich alle möglichen Zustände, die ein Prototyp erreichen kann. Jeder Zustand beschreibt eine bestimmte Situation, indem er Interaktions- und Visualisierungselementen Werte zuweist. Diese Zuweisung kann zudem an Bedingungen geknüpft sein.

\textbf{Übergänge.} Übergänge bestimmen, wie sich ein virtueller Prototyp von einem Zustand in einen anderen verändert. Diese Übergänge können durch Zeitüberschreitungen oder Ereignisse (Events) ausgelöst werden. Zusätzlich kann ein Übergang an eine zu erfüllende Bedingung (\textit{Guard}) geknüpft sein.


\section{Large Language Models und agentische Systeme}\label{sec:llms_and_agents}

Large Language Models (LLMs) sind tiefe neuronale Netzwerke, welche die Verarbeitung von natürlicher Sprache ermöglichen~\cite{lazerSurveyAgenticAI2026}.
Sie enthalten typischerweise Milliarden von Parametern und sind in der Lage, komplexe Aufgaben zu verstehen und zu lösen oder in Teilaufgaben zu zerlegen~\cite{bommasaniOpportunitiesRisksFoundation2022}.
Ihre Skalierbarkeit ist ein zentrales Merkmal von LLMs, denn sie zeigen mit zunehmender Modellgröße und Datenmenge qualitative Leistungssteigerungen und sogenannte emergente Fähigkeiten, wie \zB \textit{In-Context Learning}.
Während solche LLMs als Modelle zur passiven Sprachverarbeitung und -generierung verstanden werden können, besitzen sie ohne zusätzliche Werkzeuge keine eigene Handlungsfähigkeit.
Hier greift das Konzept eines LLM-basierten Agenten weiter.
Ein Agent ist eine autonome Einheit, welche in einer Umgebung agiert, Entscheidungen trifft und auch mit anderen Agenten interagieren kann~\cite{shohamAgentorientedProgramming1993}.
LLM-basierte Agenten verbinden diese klassische Idee mit den Fähigkeiten von LLMs.
Ein LLM kann hier als kognitive Kernkomponente verstanden werden und verleiht dem Agenten die Fähigkeit eigenständig und unabhängig zu analysieren, zu planen sowie Probleme zu lösen~\cite{xiRisePotentialLarge2023}.
So entsteht ein System, welches nicht nur natürliche Sprache verarbeiten kann, sondern eine eigene zielgerichtete Handlungsfähigkeit besitzt und in einer Umgebung agieren kann.
Ebenso kann es zusätzliche Komponenten wie Sensorik, API-Aufrufe oder Werkzeuge nutzen~\cite{xiRisePotentialLarge2023}.
Die universelle Schnittstelle für solche agentischen Systeme, über die formuliert, geplant und koordiniert wird, ist die natürliche Sprache~\cite{shenHuggingGPTSolvingAI2023}.

\subsection*{Multimodale LLMs}

Multimodale Large Language Models (MLLMs) sind Modelle, welche im Gegensatz zu reinen LLMs nicht nur Text, sondern zudem Informationen wie Bild, Audio oder Video empfangen (Input), verarbeiten sowie ausgeben (\textit{Output}) können.
Eine für diese Arbeit wichtige Fähigkeit ist ihr Verstehen von 2D-Bildern und daraus Schlussfolgerungen zu ziehen.
Sie basieren auf LLMs und bestehen im Wesentlichen aus drei Hauptmodulen: einem vortrainierten Encoder (\textit{Modality-Encoder}), einem LLM sowie einem Adapter (\textit{Connector}).
Der Encoder verarbeitet die multimodalen Signale vor und der Adapter bildet die Brücke zum LLM, indem er diese Signale für das LLM verständlich umwandelt.
Dieser Schritt ist nötig, da LLMs (wie in~\ref{sec:llms_and_agents} gezeigt) ausschließlich Text verarbeiten können. Somit sind MLLMs in der Lage, Text sowie weitere multimodale Inhalte gemeinsam zu verstehen und darauf zu antworten~\cite{yinSurveyMultimodalLarge2024, caffagniRevolutionMultimodalLarge2024}.
Im Vergleich zu früheren, ähnlichen Arbeiten basieren MLLMs auf LLMs mit Milliarden Parametern sowie neuen Trainingsparadigmen, wie dem \enquote{Multimodal Instruction Tuning}~\cite{liuVisualInstructionTuning}.
Damit sind eine Vielzahl von Möglichkeiten gegeben, welche mit reinen LLMs nicht oder nur eingeschränkt realisierbar wären.
Solche repräsentativen Fähigkeiten sind u.a. das Verständnis von visuell geprägten Witzen oder Memes~\cite{yangMMREACTPromptingChatGPT2023,zhuMiniGPT4EnhancingVisionLanguage2023}, räumliches bzw. koordinatenbezogenes Verständnis, das Erstellen von Website-Code basierend auf handgezeichneten Entwürfen oder Kochrezepten anhand von Bildern zubereiteter Gerichte~\cite{zhuMiniGPT4EnhancingVisionLanguage2023}.

\subsection*{Verwendung in der vorliegenden Arbeit}

In dieser Arbeit werden LLM-basierte Agenten als Steuerungseinheiten in einer Pipeline zur Spezifikationsgenerierung für das Vivian Framework verwendet.
Sie übernehmen eine zentrale Rolle in der Übersetzung von natürlicher Sprache in strukturierte, maschinenlesbare Spezifikationen.
Außerdem analysieren sie 3D-Szenen, planen mögliche Interaktionen mit einem 3D-Objekt, erzeugen Spezifikationen zur Realisierung dieser Interaktionen und validieren semantische Korrektheit der erzeugten Artefakte.





